% Modelo de negocio, primera entrega. Máximo 5 páginas, contando el canvas.
% Sustituir cada texto gris (\pendiente{...}) por el contenido del equipo.
% No cambiar los títulos ni el orden de las secciones.
% El canvas se edita en canvas.tex y se inserta aquí con \input. La tabla de
% hipótesis se genera sola a partir de hipotesis.csv: se actualiza el CSV, no
% este archivo. Los resultados de las pruebas van en validacion.tex.
% Compilar dos veces con: pdflatex modelo_negocio.tex
\documentclass[11pt]{article}
\usepackage[utf8]{inputenc}
\usepackage[T1]{fontenc}
\usepackage[spanish,es-noindentfirst,es-nodecimaldot,es-tabla]{babel}
\usepackage[letterpaper,margin=2.2cm,top=1.6cm]{geometry}
\usepackage{xcolor,graphicx,tikz,amssymb}
\usepackage{booktabs,tabularx,array,titling,enumitem,float,pdflscape,afterpage,datatool}
\usepackage[font=small,labelfont=bf,skip=4pt]{caption}
\usepackage[hidelinks]{hyperref}

\definecolor{unamAzul}{HTML}{003D79}
\definecolor{unamOro}{HTML}{D59F0F}
\newcommand{\pendiente}[1]{\textcolor{gray}{\emph{#1}}}
\newcolumntype{L}[1]{>{\raggedright\arraybackslash}p{#1}}
\setlist{nosep,leftmargin=*}
\setlength{\parskip}{4pt}
\setlength{\parindent}{0pt}

% ---------- Hipótesis ----------
% csv-content=literal permite %, &, $ y _ en los textos del CSV. Las claves
% siguen el orden de columnas de la plantilla, que no se cambia.
\DTLread[name=hip,format=csv,csv-content=literal,
  keys={id,bloque,hipotesis,prueba,criterio,estado,evidencia,cambio}]{hipotesis.csv}

% ---------- Encabezado institucional ----------
\pretitle{%
  \begin{center}\small
    Universidad Nacional Autónoma de México\\
    Posgrado en Ciencia e Ingeniería de la Computación\\[2pt]
    Inteligencia Artificial, semestre 2027-1
  \end{center}
  \vspace{-4pt}{\color{unamOro}\rule{\textwidth}{1.2pt}}\par\vspace{14pt}
  \begin{center}{\small\scshape Modelo de negocio}\\[6pt]\LARGE\bfseries\color{unamAzul}}
\posttitle{\par\end{center}\vspace{-4pt}}
\preauthor{\begin{center}\large\begin{tabular}[t]{c}}
\newcommand{\otrorenglon}{\end{tabular}\par\vspace{8pt}\begin{tabular}[t]{c}}
\postauthor{\end{tabular}\par\end{center}}
\predate{\begin{center}\small}
\postdate{\par\end{center}}

\title{GoIA: Navegando la ciudad de forma segura}
\author{%
  Belmont Hernández Héctor Antonio\\ {\small\texttt{habelmont2@gmail.com}}
  \and
  Hernández González Luis Ángel\\ {\small\texttt{luiss.angel92@gmail.com}}
  \otrorenglon
  Núñez Laces Yeni Yuzelmi\\ {\small\texttt{yenilaces@gmail.com}}
  \and
  Villa Aguilar Gerardo Erick\\ {\small\texttt{gerardoerickv@gmail.com}}
}
\date{Versión \pendiente{1}, \today}

\begin{document}
\maketitle
\thispagestyle{empty}

% ---------------------------------------------------------------------------
% EL MODELO EN UN PÁRRAFO. Qué se vende, a quién, qué dolor elimina, cómo llega
% al cliente y cómo se cobra. Quien lea solo este párrafo entiende el negocio.
% ---------------------------------------------------------------------------
\section{El modelo en un párrafo}
\label{sec:resumen}

GoIA es una aplicación gratuita que sugiere opciones de rutas para llegar de un punto A a un punto B considerando la seguridad de las zonas del trayecto, a comparación de solo entregar la ruta más rápida. La app está dirigida a las personas que se mueven por zonas que no
conocen, por ejemplo porque acaban de mudarse a la ciudad o la visitan, o para locales que no recorren ciertas zonas de su propia ciudad cotidianamente. El riesgo base de cada calle se calcula con los robos a transeúnte denunciados ante la Fiscalía de la CDMX, y los reportes de los usuarios lo actualizan. 

Quien
paga son los negocios ubicados sobre las rutas peatonales, que obtendrían una ubicación
destacada en el mapa, con dos de sus productos o servicios a la vista con foto y precio, para atraer a las personas que pasan a pie como forma de publicidad. La difusión se haría por medio de redes sociales principalmente. Como propuesta inicial, se cobraría 100 pesos por cada 1,000
impresiones, después de una prueba gratuita de 2,000 impresiones. La hipótesis más riesgosa es que
los usuarios perciban como más segura la ruta de GoIA que la de su planificador habitual (H3),
porque sin usuarios no hay impresiones que vender.


% ---------------------------------------------------------------------------
% CANVAS. Se inserta canvas.tex en una página horizontal, la que sigue a esta.
% El canvas se edita en canvas.tex, no aquí.
% ---------------------------------------------------------------------------
\afterpage{%
  \newgeometry{margin=1cm}%
  \begin{landscape}\thispagestyle{empty}% Business Model Canvas. Una página horizontal dentro de modelo_negocio.tex, que
% lo inserta con % Business Model Canvas. Una página horizontal dentro de modelo_negocio.tex, que
% lo inserta con % Business Model Canvas. Una página horizontal dentro de modelo_negocio.tex, que
% lo inserta con \input{canvas}. Este archivo no se compila solo.
% Llenar los nueve bloques. Cada afirmación termina con el identificador de la
% hipótesis que la respalda en hipotesis.csv, por ejemplo (H1).
% Marcar las hipótesis validadas con \val{H1} y las invalidadas con \inval{H2}.
\definecolor{navy}{HTML}{1F3467}
\providecommand{\val}[1]{\textcolor{green!50!black}{\textbf{(#1 \checkmark)}}}
\providecommand{\inval}[1]{\textcolor{red!70!black}{\textbf{(#1 $\times$)}}}

% ---------- Contenido de los bloques ----------
\newcommand{\Socios}{\begin{itemize}
  \item Una dependencia del gobierno de la CDMX acepta entregar datos de incidentes con actualización más frecuente que la de sus datos abiertos. (H9)
\end{itemize}}
\newcommand{\Actividades}{\begin{itemize}
  \item Una comunidad de 1,000 usuarios activos genera reportes suficientes para mantener actualizado el mapa de una alcaldía. (H10)
\end{itemize}}
\newcommand{\Recursos}{\begin{itemize}
  \item El riesgo base de cada calle viene de los robos a transeúnte denunciados ante la Fiscalía de la CDMX, y los usuarios lo actualizan con reportes colaborativos de incidentes o zonas inseguras. \val{H8}
\end{itemize}}
\newcommand{\Propuesta}{\begin{itemize}
  \item Los usuarios perciben como más segura la ruta que sugiere la aplicación que la de su planificador habitual. \val{H3}
  \item La ubicación destacada, que muestra sola en el mapa dos productos o servicios del negocio con foto y precio, aumenta al menos 10\% las visitas de clientes a un negocio. \val{H4}
\end{itemize}}
\newcommand{\Relacion}{\begin{itemize}
  \item Los dueños de negocios aceptan una prueba gratuita de 2,000 impresiones con reporte de visitas antes de pagar. \val{H7}
\end{itemize}}
\newcommand{\Canales}{\begin{itemize}
  \item Los dueños de negocios conocen nuevas herramientas de publicidad mediante publicidad masiva en redes sociales, radio o anuncios en vía pública. \val{H5}
  \item Al menos 3\% de las visitas a la landing page se registran para la versión de prueba. \val{H6}
\end{itemize}}
\newcommand{\Segmentos}{\begin{itemize}
  \item Las personas que se desplazan a pie o en transporte público en la CDMX cambian su ruta por seguridad aunque su planificador habitual les sugiera otra más corta. \val{H1}
  \item Los negocios de comercio y servicios ubicados sobre las rutas peatonales de la CDMX buscan más visitas de personas que pasan a pie. \val{H2}
\end{itemize}}
\newcommand{\Costos}{\begin{itemize}
  \item El servidor y el servicio de mapas para 1,000 usuarios activos cuestan menos de 3,000 pesos al mes. \val{H12}
\end{itemize}}
\newcommand{\Ingresos}{\begin{itemize}
  \item Los negocios pagarían 100 pesos por cada 1,000 impresiones como ubicación destacada. \val{H11}
\end{itemize}}

% ---------- Dibujo. No editar debajo de esta línea ----------
\begingroup
\fontsize{9}{11}\selectfont
\setlist{nosep,leftmargin=1em}
\noindent{\Large\bfseries\color{navy} Business Model Canvas}\hfill
{\large GoIA}\hfill 2.0

\vspace{4pt}
\noindent
\resizebox{\linewidth}{!}{%
\begin{tikzpicture}[x=1cm,y=1cm]
  \def\W{25.8}\def\H{16.8}\def\h{5.0}
  \tikzset{caja/.style={draw=navy,line width=0.8pt},
           tit/.style={anchor=north west,font=\bfseries\color{navy}},
           cue/.style={anchor=north west,text width=#1cm,align=left}}
  % Retícula: 5 columnas arriba, 2 abajo
  \def\c{5.16}
  \draw[caja] (0,0) rectangle (\W,\H);
  \draw[caja] (0,\h) -- (\W,\h);
  \foreach \i in {1,2,3,4} \draw[caja] (\i*\c,\h) -- (\i*\c,\H);
  \draw[caja] (\W/2,0) -- (\W/2,\h);
  \pgfmathsetmacro\m{(\h+\H)/2}
  \draw[caja] (\c,\m) -- (2*\c,\m);
  \draw[caja] (3*\c,\m) -- (4*\c,\m);
  % Títulos y contenido
  \node[tit] at (0.1,\H-0.1) {Socios clave};
  \node[cue=4.9] at (0.1,\H-0.7) {\Socios};
  \node[tit] at (\c+0.1,\H-0.1) {Actividades clave};
  \node[cue=4.9] at (\c+0.1,\H-0.7) {\Actividades};
  \node[tit] at (\c+0.1,\m-0.1) {Recursos clave};
  \node[cue=4.9] at (\c+0.1,\m-0.7) {\Recursos};
  \node[tit] at (2*\c+0.1,\H-0.1) {Propuesta de valor};
  \node[cue=4.9] at (2*\c+0.1,\H-0.7) {\Propuesta};
  \node[tit] at (3*\c+0.1,\H-0.1) {Relación con clientes};
  \node[cue=4.9] at (3*\c+0.1,\H-0.7) {\Relacion};
  \node[tit] at (3*\c+0.1,\m-0.1) {Canales};
  \node[cue=4.9] at (3*\c+0.1,\m-0.7) {\Canales};
  \node[tit] at (4*\c+0.1,\H-0.1) {Segmentos de clientes};
  \node[cue=4.9] at (4*\c+0.1,\H-0.7) {\Segmentos};
  \node[tit] at (0.1,\h-0.1) {Estructura de costos};
  \node[cue=12.6] at (0.1,\h-0.7) {\Costos};
  \node[tit] at (\W/2+0.1,\h-0.1) {Fuentes de ingresos};
  \node[cue=12.6] at (\W/2+0.1,\h-0.7) {\Ingresos};
\end{tikzpicture}}

\vspace{2pt}
\noindent\scriptsize Basado en Osterwalder y Pigneur, \emph{Business Model Generation} (2010), CC BY-SA 3.0.
\val{H1} validada,\ \inval{H2} invalidada, (H3) pendiente. Evidencia en \texttt{hipotesis.csv} y \texttt{entrevistas.csv}.
\par\endgroup
. Este archivo no se compila solo.
% Llenar los nueve bloques. Cada afirmación termina con el identificador de la
% hipótesis que la respalda en hipotesis.csv, por ejemplo (H1).
% Marcar las hipótesis validadas con \val{H1} y las invalidadas con \inval{H2}.
\definecolor{navy}{HTML}{1F3467}
\providecommand{\val}[1]{\textcolor{green!50!black}{\textbf{(#1 \checkmark)}}}
\providecommand{\inval}[1]{\textcolor{red!70!black}{\textbf{(#1 $\times$)}}}

% ---------- Contenido de los bloques ----------
\newcommand{\Socios}{\begin{itemize}
  \item Una dependencia del gobierno de la CDMX acepta entregar datos de incidentes con actualización más frecuente que la de sus datos abiertos. (H9)
\end{itemize}}
\newcommand{\Actividades}{\begin{itemize}
  \item Una comunidad de 1,000 usuarios activos genera reportes suficientes para mantener actualizado el mapa de una alcaldía. (H10)
\end{itemize}}
\newcommand{\Recursos}{\begin{itemize}
  \item El riesgo base de cada calle viene de los robos a transeúnte denunciados ante la Fiscalía de la CDMX, y los usuarios lo actualizan con reportes colaborativos de incidentes o zonas inseguras. \val{H8}
\end{itemize}}
\newcommand{\Propuesta}{\begin{itemize}
  \item Los usuarios perciben como más segura la ruta que sugiere la aplicación que la de su planificador habitual. \val{H3}
  \item La ubicación destacada, que muestra sola en el mapa dos productos o servicios del negocio con foto y precio, aumenta al menos 10\% las visitas de clientes a un negocio. \val{H4}
\end{itemize}}
\newcommand{\Relacion}{\begin{itemize}
  \item Los dueños de negocios aceptan una prueba gratuita de 2,000 impresiones con reporte de visitas antes de pagar. \val{H7}
\end{itemize}}
\newcommand{\Canales}{\begin{itemize}
  \item Los dueños de negocios conocen nuevas herramientas de publicidad mediante publicidad masiva en redes sociales, radio o anuncios en vía pública. \val{H5}
  \item Al menos 3\% de las visitas a la landing page se registran para la versión de prueba. \val{H6}
\end{itemize}}
\newcommand{\Segmentos}{\begin{itemize}
  \item Las personas que se desplazan a pie o en transporte público en la CDMX cambian su ruta por seguridad aunque su planificador habitual les sugiera otra más corta. \val{H1}
  \item Los negocios de comercio y servicios ubicados sobre las rutas peatonales de la CDMX buscan más visitas de personas que pasan a pie. \val{H2}
\end{itemize}}
\newcommand{\Costos}{\begin{itemize}
  \item El servidor y el servicio de mapas para 1,000 usuarios activos cuestan menos de 3,000 pesos al mes. \val{H12}
\end{itemize}}
\newcommand{\Ingresos}{\begin{itemize}
  \item Los negocios pagarían 100 pesos por cada 1,000 impresiones como ubicación destacada. \val{H11}
\end{itemize}}

% ---------- Dibujo. No editar debajo de esta línea ----------
\begingroup
\fontsize{9}{11}\selectfont
\setlist{nosep,leftmargin=1em}
\noindent{\Large\bfseries\color{navy} Business Model Canvas}\hfill
{\large GoIA}\hfill 2.0

\vspace{4pt}
\noindent
\resizebox{\linewidth}{!}{%
\begin{tikzpicture}[x=1cm,y=1cm]
  \def\W{25.8}\def\H{16.8}\def\h{5.0}
  \tikzset{caja/.style={draw=navy,line width=0.8pt},
           tit/.style={anchor=north west,font=\bfseries\color{navy}},
           cue/.style={anchor=north west,text width=#1cm,align=left}}
  % Retícula: 5 columnas arriba, 2 abajo
  \def\c{5.16}
  \draw[caja] (0,0) rectangle (\W,\H);
  \draw[caja] (0,\h) -- (\W,\h);
  \foreach \i in {1,2,3,4} \draw[caja] (\i*\c,\h) -- (\i*\c,\H);
  \draw[caja] (\W/2,0) -- (\W/2,\h);
  \pgfmathsetmacro\m{(\h+\H)/2}
  \draw[caja] (\c,\m) -- (2*\c,\m);
  \draw[caja] (3*\c,\m) -- (4*\c,\m);
  % Títulos y contenido
  \node[tit] at (0.1,\H-0.1) {Socios clave};
  \node[cue=4.9] at (0.1,\H-0.7) {\Socios};
  \node[tit] at (\c+0.1,\H-0.1) {Actividades clave};
  \node[cue=4.9] at (\c+0.1,\H-0.7) {\Actividades};
  \node[tit] at (\c+0.1,\m-0.1) {Recursos clave};
  \node[cue=4.9] at (\c+0.1,\m-0.7) {\Recursos};
  \node[tit] at (2*\c+0.1,\H-0.1) {Propuesta de valor};
  \node[cue=4.9] at (2*\c+0.1,\H-0.7) {\Propuesta};
  \node[tit] at (3*\c+0.1,\H-0.1) {Relación con clientes};
  \node[cue=4.9] at (3*\c+0.1,\H-0.7) {\Relacion};
  \node[tit] at (3*\c+0.1,\m-0.1) {Canales};
  \node[cue=4.9] at (3*\c+0.1,\m-0.7) {\Canales};
  \node[tit] at (4*\c+0.1,\H-0.1) {Segmentos de clientes};
  \node[cue=4.9] at (4*\c+0.1,\H-0.7) {\Segmentos};
  \node[tit] at (0.1,\h-0.1) {Estructura de costos};
  \node[cue=12.6] at (0.1,\h-0.7) {\Costos};
  \node[tit] at (\W/2+0.1,\h-0.1) {Fuentes de ingresos};
  \node[cue=12.6] at (\W/2+0.1,\h-0.7) {\Ingresos};
\end{tikzpicture}}

\vspace{2pt}
\noindent\scriptsize Basado en Osterwalder y Pigneur, \emph{Business Model Generation} (2010), CC BY-SA 3.0.
\val{H1} validada,\ \inval{H2} invalidada, (H3) pendiente. Evidencia en \texttt{hipotesis.csv} y \texttt{entrevistas.csv}.
\par\endgroup
. Este archivo no se compila solo.
% Llenar los nueve bloques. Cada afirmación termina con el identificador de la
% hipótesis que la respalda en hipotesis.csv, por ejemplo (H1).
% Marcar las hipótesis validadas con \val{H1} y las invalidadas con \inval{H2}.
\definecolor{navy}{HTML}{1F3467}
\providecommand{\val}[1]{\textcolor{green!50!black}{\textbf{(#1 \checkmark)}}}
\providecommand{\inval}[1]{\textcolor{red!70!black}{\textbf{(#1 $\times$)}}}

% ---------- Contenido de los bloques ----------
\newcommand{\Socios}{\begin{itemize}
  \item Una dependencia del gobierno de la CDMX acepta entregar datos de incidentes con actualización más frecuente que la de sus datos abiertos. (H9)
\end{itemize}}
\newcommand{\Actividades}{\begin{itemize}
  \item Una comunidad de 1,000 usuarios activos genera reportes suficientes para mantener actualizado el mapa de una alcaldía. (H10)
\end{itemize}}
\newcommand{\Recursos}{\begin{itemize}
  \item El riesgo base de cada calle viene de los robos a transeúnte denunciados ante la Fiscalía de la CDMX, y los usuarios lo actualizan con reportes colaborativos de incidentes o zonas inseguras. \val{H8}
\end{itemize}}
\newcommand{\Propuesta}{\begin{itemize}
  \item Los usuarios perciben como más segura la ruta que sugiere la aplicación que la de su planificador habitual. \val{H3}
  \item La ubicación destacada, que muestra sola en el mapa dos productos o servicios del negocio con foto y precio, aumenta al menos 10\% las visitas de clientes a un negocio. \val{H4}
\end{itemize}}
\newcommand{\Relacion}{\begin{itemize}
  \item Los dueños de negocios aceptan una prueba gratuita de 2,000 impresiones con reporte de visitas antes de pagar. \val{H7}
\end{itemize}}
\newcommand{\Canales}{\begin{itemize}
  \item Los dueños de negocios conocen nuevas herramientas de publicidad mediante publicidad masiva en redes sociales, radio o anuncios en vía pública. \val{H5}
  \item Al menos 3\% de las visitas a la landing page se registran para la versión de prueba. \val{H6}
\end{itemize}}
\newcommand{\Segmentos}{\begin{itemize}
  \item Las personas que se desplazan a pie o en transporte público en la CDMX cambian su ruta por seguridad aunque su planificador habitual les sugiera otra más corta. \val{H1}
  \item Los negocios de comercio y servicios ubicados sobre las rutas peatonales de la CDMX buscan más visitas de personas que pasan a pie. \val{H2}
\end{itemize}}
\newcommand{\Costos}{\begin{itemize}
  \item El servidor y el servicio de mapas para 1,000 usuarios activos cuestan menos de 3,000 pesos al mes. \val{H12}
\end{itemize}}
\newcommand{\Ingresos}{\begin{itemize}
  \item Los negocios pagarían 100 pesos por cada 1,000 impresiones como ubicación destacada. \val{H11}
\end{itemize}}

% ---------- Dibujo. No editar debajo de esta línea ----------
\begingroup
\fontsize{9}{11}\selectfont
\setlist{nosep,leftmargin=1em}
\noindent{\Large\bfseries\color{navy} Business Model Canvas}\hfill
{\large GoIA}\hfill 2.0

\vspace{4pt}
\noindent
\resizebox{\linewidth}{!}{%
\begin{tikzpicture}[x=1cm,y=1cm]
  \def\W{25.8}\def\H{16.8}\def\h{5.0}
  \tikzset{caja/.style={draw=navy,line width=0.8pt},
           tit/.style={anchor=north west,font=\bfseries\color{navy}},
           cue/.style={anchor=north west,text width=#1cm,align=left}}
  % Retícula: 5 columnas arriba, 2 abajo
  \def\c{5.16}
  \draw[caja] (0,0) rectangle (\W,\H);
  \draw[caja] (0,\h) -- (\W,\h);
  \foreach \i in {1,2,3,4} \draw[caja] (\i*\c,\h) -- (\i*\c,\H);
  \draw[caja] (\W/2,0) -- (\W/2,\h);
  \pgfmathsetmacro\m{(\h+\H)/2}
  \draw[caja] (\c,\m) -- (2*\c,\m);
  \draw[caja] (3*\c,\m) -- (4*\c,\m);
  % Títulos y contenido
  \node[tit] at (0.1,\H-0.1) {Socios clave};
  \node[cue=4.9] at (0.1,\H-0.7) {\Socios};
  \node[tit] at (\c+0.1,\H-0.1) {Actividades clave};
  \node[cue=4.9] at (\c+0.1,\H-0.7) {\Actividades};
  \node[tit] at (\c+0.1,\m-0.1) {Recursos clave};
  \node[cue=4.9] at (\c+0.1,\m-0.7) {\Recursos};
  \node[tit] at (2*\c+0.1,\H-0.1) {Propuesta de valor};
  \node[cue=4.9] at (2*\c+0.1,\H-0.7) {\Propuesta};
  \node[tit] at (3*\c+0.1,\H-0.1) {Relación con clientes};
  \node[cue=4.9] at (3*\c+0.1,\H-0.7) {\Relacion};
  \node[tit] at (3*\c+0.1,\m-0.1) {Canales};
  \node[cue=4.9] at (3*\c+0.1,\m-0.7) {\Canales};
  \node[tit] at (4*\c+0.1,\H-0.1) {Segmentos de clientes};
  \node[cue=4.9] at (4*\c+0.1,\H-0.7) {\Segmentos};
  \node[tit] at (0.1,\h-0.1) {Estructura de costos};
  \node[cue=12.6] at (0.1,\h-0.7) {\Costos};
  \node[tit] at (\W/2+0.1,\h-0.1) {Fuentes de ingresos};
  \node[cue=12.6] at (\W/2+0.1,\h-0.7) {\Ingresos};
\end{tikzpicture}}

\vspace{2pt}
\noindent\scriptsize Basado en Osterwalder y Pigneur, \emph{Business Model Generation} (2010), CC BY-SA 3.0.
\val{H1} validada,\ \inval{H2} invalidada, (H3) pendiente. Evidencia en \texttt{hipotesis.csv} y \texttt{entrevistas.csv}.
\par\endgroup
\end{landscape}%
  \restoregeometry}

% ---------------------------------------------------------------------------
% CLIENTE Y PROPUESTA DE VALOR. La persona es un personaje ficticio construido
% con hechos de clientes reales, no con suposiciones. Cada celda cita las
% entrevistas (E01) o las hipótesis (H1) que la respaldan. Si quien paga y quien
% usa son la misma persona, borrar la columna Usuario.
% ---------------------------------------------------------------------------
\section{Cliente y propuesta de valor}
\label{sec:cliente}

El segmento principal son las personas que se desplazan a pie o en transporte público por zonas de
la CDMX que no conocen, como quienes acaban de mudarse o la visitan. Se eligió sobre personas locales porque ellos saben qué calles evitar. 
En contraste, quien paga es un segmento distinto, los negocios sobre las rutas peatonales, porque se
benefician directamente de las personas que la ruta lleva frente a su local (H2). La
Tabla~\ref{tab:persona} describe perfiles ficticios de los arqueotipos que pagarían y usarían el producto.

\begin{table}[H]
  \caption{Perfil de la persona que paga y de la que usa el producto. Cada celda cita la
    entrevista o la hipótesis que la respalda.}
  \label{tab:persona}
  \small
  \renewcommand{\arraystretch}{1.25}
  \begin{tabularx}{\textwidth}{@{}>{\bfseries\raggedright}p{0.17\textwidth}
                                  >{\raggedright\arraybackslash}X>{\raggedright\arraybackslash}X@{}}
    \toprule
    & Cliente que paga & Usuario \\
    \midrule
    Nombre y cargo & Marta, dueña de una cafetería. & Daniela, estudiante de posgrado que acaba
                     de mudarse a la CDMX. \\
    Organización & Cafetería con dos empleados en una calle de paso de la
                   CDMX. (H2) & Ninguna. Se traslada por su cuenta. \\
    Un día típico & Atiende el local. Sus ventas dependen de cuántas personas pasan frente a
                    él. (H2) & Camina del metro a la universidad y por colonias que todavía no
                    conoce. Google Maps le sugiere la ruta más rápida, que puede cruzar una calle
                    insegura sin que ella lo sepa. (H1) \\
    Objetivos & Atraer más clientes de paso. (H2) & Moverse por la ciudad como alguien que
                vive ahí, sin exponerse a un asalto. (H1) \\
    Dolores & Paga publicidad local o en internet sin saber cuántos clientes le trae.
              (H2, H11) & No sabe qué zonas evitar. Un asalto le costaría la cartera, el celular
              y hasta el vaciado de sus cuentas. (H1) \\
    Ganancias & Ver cuántos clientes llegaron por la aplicación. (H4, H7) & Una ruta que
                perciba como segura, sin pasar de un tiempo de llegada razonable, y poder
                compararla con la más rápida. (H3) \\
    Solución actual & Volantes, redes sociales y anuncios pagados en internet. (H2) & Google
                      Maps, que solo minimiza el tiempo, y preguntar a compañeros o
                      vecinos. (H1) \\
    Cómo compra & Decide y paga ella, después de una prueba gratuita. (H7, H11) & No paga. Usa
                  la aplicación gratuita. \\
    Dónde se informa & Redes sociales, radio y anuncios en vía pública. (H5) & LinkedIn, grupos
                       de estudiantes, grupos de la CDMX y videos cortos. (H6) \\
    \bottomrule
  \end{tabularx}
\end{table}


% ---------------------------------------------------------------------------
% RELACIÓN CON CLIENTES Y CANALES. Los canales cubren las cinco fases del
% recorrido del cliente. Un canal concreto dice a quién llega y con qué mensaje:
% "redes sociales" no basta, "anuncios en LinkedIn a jefes de despacho" sí. La
% relación dice cómo se consigue, se conserva y se hace crecer a cada cliente.
% El costo de cada canal y de cada relación alimenta el CAC del modelo financiero.
% ---------------------------------------------------------------------------
\section{Relación con clientes y canales}
\label{sec:canales}

El principal canal para adquirir usuarios será una landing page promocionada en redes sociales, grupos universitarios y comunidades locales de la Ciudad de México, ya que permite llegar directamente a personas que realizan desplazamientos frecuentes y que podrían estar interesadas en rutas más seguras. Posteriormente, la misma aplicación fomentará el crecimiento mediante reportes colaborativos y recomendaciones entre usuarios. La Tabla~\ref{tab:canales} describe los canales en
cada fase, y la Tabla~\ref{tab:relacion} describe la relación con el cliente.

\begin{table}[H]
	\caption{Canales en las cinco fases del recorrido del cliente. Directo significa que el
		equipo opera el canal, e indirecto que lo opera un socio.}
	\label{tab:canales}
	\small
	\renewcommand{\arraystretch}{1.2}
	\begin{tabularx}{\textwidth}{@{}>{\bfseries}p{0.15\textwidth}Xp{0.11\textwidth}p{0.16\textwidth}c@{}}
		\toprule
		Fase & Canal & Tipo & Costo estimado & Hipótesis \\
		\midrule
		Conocimiento & Publicaciones en redes sociales, grupos universitarios y comunidades locales & Directo& $\$500$ MXN inicialmente& H6\\
		Evaluación & Landing page con demostración de rutas seguras y comparación con Google Maps & Directo & $\$ 0 - \$300$ MXN mensuales & H1, H3  \\
		Compra & Registro gratuito mediante correo electrónico o cuenta Google  & Directo &  Sin costo para el usuario& H6 \\
		Entrega & Aplicación web accesible desde navegador móvil o computadora & Directo & Incluido en costos de infraestructura & H3 \\
		Posventa & Notificaciones, formulario de retroalimentación y sistema de reportes & Directo & Sin costo adicional& H8 \\
		\bottomrule
	\end{tabularx}
\end{table}

\begin{table}[H]
	\caption{Relación con el cliente en tres etapas. Los tipos son asistencia personal,
		asistencia dedicada, autoservicio, servicio automatizado, comunidad y cocreación.}
	\label{tab:relacion}
	\small
	\renewcommand{\arraystretch}{1.2}
	\begin{tabularx}{\textwidth}{@{}>{\bfseries}p{0.16\textwidth}p{0.2\textwidth}Xc@{}}
		\toprule
		Etapa & Tipo & Cómo se hace y cuánto cuesta & Hipótesis \\
		\midrule
		Personas que caminan o utilizan transporte público & Autoservicio & El usuario consulta rutas y niveles de riesgo sin intervención humana& H1 \\
		Personas preocupadas por la seguridad de sus trayectos &Comunidad &Se recopila retroalimentación para mejorar los criterios de riesgo y las recomendaciones  & H8 \\
		Negocios locales interesados en visibilidad & Asistencia automatizada & Contratación de ubicaciones destacadas dentro de la plataforma & H4, H11\\
		\bottomrule
	\end{tabularx}
\end{table}


% ---------------------------------------------------------------------------
% COMPETENCIA. Todo problema real tiene una solución actual, aunque sea una hoja
% de cálculo. Incluir esa solución y al menos dos competidores directos. Los dos
% atributos son los que más le importan al cliente según las entrevistas.
% ---------------------------------------------------------------------------
\section{Competencia}
\label{sec:competencia}

La Tabla~\ref{tab:competencia} compara las opciones del cliente en precio, facilidad y seguridad, los tres atributos que más mencionaron las personas entrevistadas.

\begin{table}[H]
  \caption{Opciones actuales del cliente. Calificación de 1 (bajo) a 5 (alto) en los dos
    atributos que más le importan, y precio con su fuente.}
  \label{tab:competencia}
  \small
  \renewcommand{\arraystretch}{1.2}
  \begin{tabularx}{\textwidth}{@{}Xp{0.1\textwidth}cccp{0.2\textwidth}@{}}
    \toprule
    Opción & Tipo & Precio Accesible & Facilidad de Uso & Seguridad & Precio y fuente \\
    \midrule
    Google Maps & Indirecto & 5 &  5& 2 & Gratuito Versión Básica\\
    Waze  & Indirecto & 5 &  4& 3 & Gratuito Versión Básica \\
    Conocimiento Propio & Solución actual & 5 &  3& 3 & Sin costo directo, basado en experiencia y recomendaciones.\\
    \textbf{GoIA}  & Propuesta propia & 5 & 4 & 5 & Gratuito en etapa de validación.\\
    \bottomrule
  \end{tabularx}
\end{table}

La principal ventaja de GoIA es incorporar información de seguridad dentro del resultado de la búsqueda de rutas mediante el algoritmo A*. Mientras que aplicaciones como Google Maps y Waze priorizan el tiempo de recorrido y las condiciones del tránsito, GoIA considera además el nivel de riesgo asociado a cada zona de una ruta. La aplicación contempla la retroalimentación de los usuarios para reportar incidentes y zonas inseguras (H8), permitiendo mantener actualizada la evaluación de riesgo de las ruta.

% ---------------------------------------------------------------------------
% TAMAÑO DE MERCADO. De abajo hacia arriba: número de clientes por precio anual.
% Cada número lleva su fuente. Las cifras son las mismas de la hoja Supuestos del
% modelo financiero.
% ---------------------------------------------------------------------------
\section{Tamaño de mercado}
\label{sec:mercado}

La Tabla~\ref{tab:mercado} estima el mercado de abajo hacia arriba, con el número de clientes
por el precio anual. 
%\pendiente{Una oración que justifica el porcentaje que lleva del TAM al
%SAM y del SAM al SOM.}



Además, un análisis de robo a transeúnte la coloca en primer lugar entre las alcaldías de CDMX, con una tasa de 390.8 por cada 100 mil habitantes en el periodo analizado

\url{https://enfemex.com/datos/seguridad/robo-transeunte}


\begin{table}[H]
  \caption{Tamaño de mercado. TAM es todo el mercado si cada cliente posible comprara, SAM es la
    parte que alcanzan los canales del equipo y SOM es la parte que el equipo espera ganar en 24
    meses. Montos en pesos por año.}
  \label{tab:mercado}
  \small
  \renewcommand{\arraystretch}{1.2}
  \begin{tabularx}{\textwidth}{@{}lrrrX@{}}
    \toprule
    Nivel & Clientes & Precio anual & Mercado & Fuente o supuesto \\
    \midrule
    TAM & $ 63,215$ & $ \$4200$ & $ \$265,503,000$& Para el mercado consideramos un CPM de $\$$100, un supuesto de 15 impresiones mensuales por usuario. 180 impresiones anuales, lo que da un ingreso por usuario anual de $\$$18.
     \url{https://chatgpt.com/c/6af3c2b6-8bc4-83e8-a3d3-b54b74fd30f5}
    
     \\
    SAM & $12,643$ & $ \$4200$ & $\$53,100,000$ & Se estima que el 20$\%$ corresponde a negocios potencialmente interesados en publicidad geolocalizada. 
     \url{https://en.www.inegi.org.mx/contenidos/saladeprensa/boletines/2025/ce/CE_2024_Def_CdMx.pdf?} \\
    SOM & $ 505$ & $ \$4200$ & $\$2,121,000$ & Pretendemos alcanzar en el primer año el 10$\%$ del SOM por lo que tendríamos 2000 usuarios activos. \\
    \bottomrule
  \end{tabularx}
\end{table}

% ---------------------------------------------------------------------------
% HIPÓTESIS Y PRUEBAS. La tabla se genera a partir de hipotesis.csv. Cada
% afirmación del canvas tiene una hipótesis, y cada hipótesis tiene una prueba y
% un criterio de éxito fijado antes de probarla. Los resultados van en
% validacion.tex, en la entrega final.
% ---------------------------------------------------------------------------
\section{Hipótesis y pruebas}
\label{sec:hipotesis}

\pendiente{Un párrafo que dice cuáles son las dos o tres hipótesis más riesgosas, las que
harían fracasar el negocio si fueran falsas, y en qué orden se probarán.} La
Tabla~\ref{tab:hipotesis} lista las hipótesis del canvas.

\begin{table}[H]
  \caption{Hipótesis del canvas, con la prueba que la pondrá a prueba y el criterio de éxito
    fijado antes de probarla.}
  \label{tab:hipotesis}
  \footnotesize
  \renewcommand{\arraystretch}{1.2}
  \begin{tabular}{@{}lL{0.38\textwidth}L{0.22\textwidth}L{0.28\textwidth}@{}}
    \toprule
    id & Hipótesis (bloque) & Prueba & Criterio de éxito \\
    \midrule
    \DTLforeach*{hip}{\hId=id,\hBloque=bloque,\hTexto=hipotesis,\hPrueba=prueba,\hCrit=criterio}{%
      \DTLiffirstrow{}{\\}%
      \hId & \hTexto{} {\color{gray}(\hBloque)} & \hPrueba & \hCrit}\\
    \bottomrule
  \end{tabular}
\end{table}

\end{document}
